% Cadre décoratif avec TikZ
\thispagestyle{empty}
 \newgeometry{top=3cm,bottom=3cm,left=2.5cm,right=2.5cm}
    {\centering
    %\begin{tikzpicture}[remember picture, overlay]
        %\draw[line width=2pt, color=black] 
       %     ([xshift=1cm,yshift=-1cm]current page.north west) 
       %    rectangle 
       %    ([xshift=-1cm,yshift=1cm]current page.south east);
    %\end{tikzpicture}
    
    % Logo ou image
    \vspace*{1cm}
    \includegraphics[width=0.5\textwidth]{\rootpath/doc/common/figures/InrUM.png} 
    \vspace*{2cm}
    
    % Titre
    {\Huge {SW2D}}\\[.3cm]
    {\Large\bfseries \guide \\[0.5cm]}
    \rule{\linewidth}{0.5mm}
    % Sous-titre
    ~\\
    {\large Version \myversion}\\[0.5cm]
    {\large \today}\\[2cm]
    
        \vspace*{1cm}
    \includegraphics[width=0.4\textwidth]{\rootpath/doc/common/figures/logo_sw2d.png}
    \vspace*{2cm}
    
    % Auteur ou société
    \vfill
    \textsc{Main developers:}\\[0.3cm]
    {\Large \authors}\\[1cm]
    
    % Bas de page
    \vfill
    \rule{0.8\textwidth}{0.4pt}\\[0.3cm]
    {\small © 2025 Inria \& Université de Montpellier}
    \restoregeometry
    }
\pagebreak

\tableofcontents

\pagebreak

\fancyfoot[R]{SW2D \guide, version \myversion}